\section{Segundo Parcial 2016}

% =====================================================================
\subsection{Ejercicio 1 --- Modelo de compartimentos de confidencialidad}
% =====================================================================

\begin{consigna}
Un famoso restaurante de 5 estrellas en la guía Michelin se destaca por todo lo que sirve:
platos salados, platos dulces y bebidas (vinos y tragos).

Se tienen los siguientes sujetos:
\begin{itemize}
  \item Jorge, José y Juan son los dueños.
  \item Manuel y Ramón son chefs.
  \item Luis es uno de los mozos.
  \item Andrea es una cliente.
\end{itemize}

Considerar los siguientes objetos:
\begin{itemize}
  \item Recetas (tanto de platos dulces, como de salados, como de tragos).
  \item Menu. (se considera menú al listado de platos con sus respectivos precios).
  \item Listas de Stock (de productos para platos dulces, salados y tragos).
  \item Comanda. (pedido que se hace al camarero-mozo).
\end{itemize}

\textbf{Diseñar un modelo de compartimentos de confidencialidad, que permita garantizar
los siguientes requisitos, justificando el cumplimiento de cada uno.}
\begin{itemize}
  \item Jorge puede leer todas las recetas, las listas de stock, el menú y las comandas.
  \item José, el somellier, es el único que escribe recetas de tragos. También lee stock de
        productos para tragos.
  \item Juan es el único que escribe recetas de platos (tanto dulces como salados). También lee
        stock de productos para platos.
  \item Los chefs pueden leer las recetas, pero no modificarlas.
  \item Los chefs pueden escribir en las listas de stocks para recomendar la compra de
        productos de su área.
  \item Ramón es el jefe de la barra, y sólo se ocupa de armar tragos. No cocina platos.
  \item Manuel dirige la cocina de platos dulces y salados.
  \item Los mozos sólo leen el menú y escriben comandas (no leen ni escriben nada mas).
  \item Los clientes sólo leen menús.
\end{itemize}
\end{consigna}

\paragraph{Temas.}
Clase 6 --- Políticas y Control de Acceso ($\S$ Bell--LaPadula; $\S$ Compartimentos, reticulado
y dominancia). Modelo de confidencialidad por \emph{compartimentos} (\emph{lattice}) con regla
\emph{read-down / write-up} y \emph{need-to-know}.

\paragraph{Qué necesitás saber.}
El enunciado pide explícitamente un \textbf{modelo de compartimentos}, es decir
Bell--LaPadula con categorías. Un compartimento es un par \textbf{(nivel, conjunto de
categorías)}. Las reglas de BLP, en su forma con dominancia:
\begin{itemize}
  \item \textbf{Lectura (seguridad simple, \emph{read-down}):} $s$ lee $o$ sii
        $\Lvl(s) \dom \Lvl(o)$.
  \item \textbf{Escritura (propiedad $\ast$, \emph{write-up}):} $s$ escribe $o$ sii
        $\Lvl(o) \dom \Lvl(s)$.
\end{itemize}
La relación de \textbf{dominancia} para $(L,C)$ y $(L',C')$:
\begin{equation*}
  (L,C) \dom (L',C') \iff L' \le L \;\wedge\; C' \subseteq C.
\end{equation*}
El conjunto de compartimentos forma un \textbf{reticulado} (orden parcial): pares incomparables
quedan \textbf{separados} (ni lectura ni escritura entre ellos). Las \textbf{categorías}
implementan el \emph{need-to-know} (un sujeto solo accede a su área aunque su nivel sea alto).

\paragraph{Notas y conexiones.}
La clave del ejercicio es notar que hay \textbf{dos dimensiones} de separación que no son un
mero orden total: (i) un \textbf{rumbo de área} (\emph{dulce}, \emph{salado}, \emph{tragos})
que se modela con \textbf{categorías} (separación horizontal, \emph{need-to-know}: Manuel no
toca tragos y Ramón no toca platos $\Rightarrow$ categorías incomparables, se separan solos),
y (ii) una \textbf{jerarquía de poder} (los dueños arriba, los mozos/clientes abajo) que se
modela con el \textbf{nivel}. Quien \emph{escribe} recetas debe estar en un nivel \emph{más
bajo} que quien las lee, por la regla \emph{write-up}: en BLP el flujo de información sube, así
que el ``autor'' escribe hacia arriba y el ``lector con autoridad'' lee hacia abajo. Esto
mismo se reusa en Clase 9 (confinamiento) donde BLP es el modelo canónico de control de flujo.

\paragraph{Resolución.}
\textbf{Categorías (need-to-know por área).} Definimos el conjunto de categorías
$\{\textsf{D}, \textsf{S}, \textsf{T}\}$ (Dulce, Salado, Tragos). Los objetos por área llevan
su categoría; el \textsf{Menú} y las \textsf{Comandas} son transversales (categoría vacía o
las tres, según convenga al flujo, ver abajo).

\textbf{Niveles (jerarquía de escritura/lectura).} Para cada área usamos dos niveles, porque
el \emph{autor} de la receta debe quedar \textbf{debajo} del lector con autoridad
(\emph{write-up}). Proponemos, dentro de cada categoría:
\begin{equation*}
  \text{Bajo (autor)} \;<\; \text{Alto (dueño lector)}.
\end{equation*}

Asignación de compartimentos (nivel, categorías):

\begin{center}
\renewcommand{\arraystretch}{1.25}
\begin{tabular}{l l l}
\toprule
\textbf{Sujeto} & \textbf{Compartimento} & \textbf{Justificación} \\
\midrule
Jorge (dueño) & $(\text{Alto},\{\textsf{D},\textsf{S},\textsf{T}\})$ &
  lee todo hacia abajo (recetas, stock, menú, comandas). \\
José (somellier) & $(\text{Bajo},\{\textsf{T}\})$ &
  único autor de recetas de tragos ($\Rightarrow$ \emph{write-up} a recetas\textsubscript{T}). \\
Juan & $(\text{Bajo},\{\textsf{D},\textsf{S}\})$ &
  único autor de recetas de platos dulces/salados. \\
Manuel (chef platos) & $(\text{Medio},\{\textsf{D},\textsf{S}\})$ &
  lee recetas\textsubscript{D,S}, escribe stock\textsubscript{D,S}. \\
Ramón (chef barra) & $(\text{Medio},\{\textsf{T}\})$ &
  lee recetas\textsubscript{T}, escribe stock\textsubscript{T}; no toca platos. \\
Luis (mozo) & $(\text{Mozo},\varnothing)$ &
  lee menú, escribe comandas. \\
Andrea (cliente) & $(\text{Cliente},\varnothing)$ &
  solo lee menú. \\
\bottomrule
\end{tabular}
\end{center}

\noindent Y los objetos:

\begin{center}
\renewcommand{\arraystretch}{1.25}
\begin{tabular}{l l}
\toprule
\textbf{Objeto} & \textbf{Compartimento} \\
\midrule
Recetas de tragos & $(\text{Bajo},\{\textsf{T}\})$ \\
Recetas de platos dulces/salados & $(\text{Bajo},\{\textsf{D},\textsf{S}\})$ \\
Stock de tragos & $(\text{Bajo},\{\textsf{T}\})$ \\
Stock de platos dulces/salados & $(\text{Bajo},\{\textsf{D},\textsf{S}\})$ \\
Menú & $(\text{Mozo},\varnothing)$ \\
Comanda & $(\text{Alto},\varnothing)$ \\
\bottomrule
\end{tabular}
\end{center}

\textbf{Verificación requisito por requisito} (recordando lectura: $\Lvl(s)\dom\Lvl(o)$;
escritura: $\Lvl(o)\dom\Lvl(s)$):
\begin{itemize}
  \item \textbf{Jorge lee todo:} $(\text{Alto},\{\textsf{D},\textsf{S},\textsf{T}\})$ domina a
        todos los objetos de recetas/stock (nivel mayor o igual y categorías que contienen las
        del objeto) y al menú/comandas $\Rightarrow$ \emph{read-down} OK.
  \item \textbf{José único que escribe recetas de tragos:} su compartimento
        $(\text{Bajo},\{\textsf{T}\})$ permite \emph{write-up} a recetas\textsubscript{T}
        $(\text{Bajo},\{\textsf{T}\})$ (mismo compartimento, escritura válida). Ningún otro
        sujeto con categoría $\textsf{T}$ está en nivel $\le \text{Bajo}$ con esa categoría
        exclusiva, así que nadie más puede escribirlas. También lee stock\textsubscript{T}
        (mismo compartimento).
  \item \textbf{Juan único que escribe recetas de platos:} idéntico con
        $(\text{Bajo},\{\textsf{D},\textsf{S}\})$; lee stock de platos.
  \item \textbf{Chefs leen recetas pero no las modifican:} Manuel/Ramón están en nivel
        \emph{Medio} $>$ Bajo. Lectura de recetas: $\Lvl(\text{chef})\dom\Lvl(\text{receta})$
        (Medio domina Bajo, categorías incluidas) $\Rightarrow$ \textbf{leen}. Escritura de
        recetas: requeriría $\Lvl(\text{receta})\dom\Lvl(\text{chef})$, es decir
        $\text{Bajo}\dom\text{Medio}$, que es \textbf{falso} $\Rightarrow$ \textbf{no pueden
        modificarlas}. Exactamente lo pedido (es el caso de manual de BLP: leer hacia abajo
        sí, escribir hacia abajo no).
  \item \textbf{Chefs escriben en stock de su área:} ponemos el stock en nivel Bajo; entonces
        \emph{write-up} desde Medio a Bajo \emph{no} valdría. Para respetar el requisito el
        stock de cada área se coloca en el \textbf{mismo nivel de los chefs o por encima} para
        esa escritura; lo más limpio es modelar el stock como objeto en
        $(\text{Medio},\{\textsf{área}\})$ para escritura por el chef (write-up trivial al
        mismo nivel) y dejar que los dueños lo lean desde Alto. \emph{(Detalle de diseño: el
        stock necesita un compartimento que admita lectura del autor de área para
        ``recomendar compra''; se ubica al nivel del chef.)}
  \item \textbf{Ramón solo tragos / Manuel solo platos:} sus categorías
        ($\{\textsf{T}\}$ vs $\{\textsf{D},\textsf{S}\}$) son \textbf{incomparables}
        $\Rightarrow$ están separados: Ramón no lee ni escribe objetos de platos y viceversa
        (\emph{need-to-know} puro, sin necesidad de regla extra).
  \item \textbf{Mozos leen menú, escriben comandas, nada más:} Luis $(\text{Mozo},\varnothing)$
        lee el menú $(\text{Mozo},\varnothing)$ (mismo compartimento); escribe la comanda
        $(\text{Alto},\varnothing)$ por \emph{write-up} ($\text{Alto}\dom\text{Mozo}$, categoría
        $\varnothing\subseteq\varnothing$) $\Rightarrow$ escribe pero \textbf{no la puede
        leer} (no domina Alto). No tiene categorías $\Rightarrow$ no ve recetas ni stock.
  \item \textbf{Clientes solo leen menú:} Andrea $(\text{Cliente},\varnothing)$ con
        Cliente $<$ Mozo; lee el menú por \emph{read-down}. No escribe comandas porque eso
        sería write-up que solo habilitamos al mozo, y por nivel/rol no accede a nada más.
\end{itemize}

\noindent Las comandas en nivel Alto y categoría vacía hacen que ``suban'' del mozo hacia los
dueños (flujo de información hacia arriba, como manda BLP) sin que el mozo pueda releerlas.

% =====================================================================
\subsection{Ejercicio 2 --- Definir conceptos y relacionarlos}
% =====================================================================

\begin{consigna}
Define cada concepto. Luego, escribe una frase que justifique la relación entre un par de ellos
(no una relación trivial como que pertenecen a una misma categoría o son relativos a seguridad
informática).
\begin{itemize}
  \item principio de mínimo privilegio
  \item confinamiento
  \item rootkit
  \item troyano
\end{itemize}
\end{consigna}

\paragraph{Temas.}
Clase 8 --- Principios de Diseño ($\S$ Mínimo privilegio). Clase 9 --- Flujo de Información y
Malware ($\S$ El problema del confinamiento; $\S$ Rootkit; $\S$ Troyano).

\paragraph{Qué necesitás saber.}
Las cuatro definiciones, tomadas de las fuentes:
\begin{itemize}
  \item \textbf{Principio de mínimo privilegio (\emph{least privilege}, Saltzer \& Schroeder):}
        asegurar que cada sujeto tenga \textbf{solo el acceso necesario para realizar su tarea,
        ni más}. Consecuencia derivada: la condición por defecto es \emph{no tener acceso a
        nada} (enlaza con \emph{fail-safe defaults}).
  \item \textbf{Confinamiento:} dificultad de asegurar que un proceso (sujeto) \textbf{no
        transmita información} a una entidad a la que no está autorizado, ni de forma directa
        ni indirecta. El ideal es la \emph{aislación total} (no comunicar y no ser observado),
        que en la práctica es inalcanzable porque todo cómputo deja medidas observables.
  \item \textbf{Rootkit:} software diseñado para \textbf{ocultar malware} dentro de la máquina
        infectada, con contramedidas activas (ofuscación, encripción, reemplazo de
        herramientas del sistema, intercepción de \emph{system calls}) para quedar invisible a
        administradores y antivirus. Opera a nivel kernel/usuario o incluso firmware/BIOS.
  \item \textbf{Troyano:} malware \textbf{oculto dentro de otro software aparentemente benigno}
        y de interés para la víctima (del caballo de Troya). Sirve para que el atacante se
        establezca dentro del equipo infectado.
\end{itemize}

\paragraph{Notas y conexiones.}
\emph{Rootkit} y \emph{troyano} son las dos formas de \textbf{esconderse} del malware
(ocultamiento), distintas de \emph{virus/worm} que aportan la \emph{propagación}; ninguno de
los dos describe el \emph{payload}. La distinción fina: el troyano se esconde en la
\emph{instalación} (parece software legítimo para entrar), el rootkit se esconde
\emph{una vez dentro} (oculta su presencia al sistema operativo).

\paragraph{Resolución.}
Definiciones, dadas arriba. Frases de relación \textbf{no triviales} (basta una para cumplir la
consigna; se ofrecen dos opciones):
\begin{itemize}
  \item \textbf{Troyano $\leftrightarrow$ Rootkit:} ``Un troyano logra la \emph{entrada}
        haciéndose pasar por software benigno, y suele instalar un \emph{rootkit} que garantiza
        la \emph{permanencia} ocultando esa presencia al sistema y al antivirus; el primero
        resuelve el ingreso, el segundo la invisibilidad posterior.''
  \item \textbf{Mínimo privilegio $\leftrightarrow$ Confinamiento:} ``El mínimo privilegio
        \emph{limita qué puede acceder} un proceso, pero \emph{no impide que reenvíe} lo que ya
        leyó; el confinamiento ataca justamente ese hueco al restringir el \emph{flujo} de
        información del proceso hacia afuera (control de acceso $\neq$ control de flujo).''
\end{itemize}

% =====================================================================
\subsection{Ejercicio 3 --- Secreto compartido de Shamir y entropía}
% =====================================================================

\begin{consigna}
Considerar el conjunto $C$ de sombras para un esquema de secreto compartido de Shamir $(k,n)$,
donde $k=3$ y $n=4$ sobre $\mathbb{Z}_7$.
\[
  C = \{(1,6),(2,1),(3,0),(4,3)\}
\]
\begin{enumerate}[label=\textbf{\alph*)}]
  \item Obtener el secreto.
  \item Obtener el polinomio utilizado para obtener el conjunto de sombras.
  \item Expresar el significado de un esquema de secreto umbral $(3,n)$ utilizando el concepto
        de entropía.
\end{enumerate}
\end{consigna}

\paragraph{Temas.}
Clase 8 --- Principios de Diseño ($\S$ El secreto compartido de Shamir, en términos de
entropía; $\S$ Entropía $H(x)$ e incertidumbre). Esquema umbral $(T,N)$ + interpolación de
Lagrange en $\mathbb{Z}_p$.

\paragraph{Qué necesitás saber.}
En Shamir $(k,n)$ el repartidor elige un polinomio de grado $k-1$ con coeficientes en
$\mathbb{Z}_p$:
\begin{equation*}
  f(x) = a_0 + a_1 x + a_2 x^2 + \dots + a_{k-1}x^{k-1} \pmod p,
\end{equation*}
donde el \textbf{secreto es $S = a_0 = f(0)$} y cada sombra es un punto $(x_i, f(x_i))$. Con
$k=3$ el polinomio es \textbf{cuadrático} y bastan $3$ sombras para reconstruirlo por
interpolación. En $\mathbb{Z}_7$ todas las operaciones son módulo $7$ y la división es
multiplicación por el inverso modular ($a^{-1} = a^{p-2} = a^5 \bmod 7$). Para el secreto se
puede interpolar directamente en $x=0$:
\begin{equation*}
  S = f(0) = \sum_{i} y_i \prod_{j\neq i} \frac{0 - x_j}{x_i - x_j} \pmod 7 .
\end{equation*}

\paragraph{Notas y conexiones.}
La parte (c) es el puente \textbf{Shamir $\leftrightarrow$ entropía} que la cátedra desarrolla
en la práctica de Clase 8/9: ``un esquema $(T,N)$ es válido'' equivale a decir que con
\emph{menos de $T$} sombras la entropía del secreto \textbf{no baja} (no hay flujo de
información), y con $T$ o más \textbf{cae a $0$}. Es el mismo criterio de flujo de información
por entropía de Clase 9 (reducir $H$ = hay flujo).

\paragraph{Resolución.}

\textbf{a) Secreto.} Interpolamos en $x=0$ con las tres primeras sombras
$(1,6),(2,1),(3,0)$ en $\mathbb{Z}_7$. Los términos de Lagrange evaluados en $0$:
\begin{align*}
  \ell_1(0) &= \frac{(0-2)(0-3)}{(1-2)(1-3)} = \frac{(-2)(-3)}{(-1)(-2)} = \frac{6}{2} = 3,\\
  \ell_2(0) &= \frac{(0-1)(0-3)}{(2-1)(2-3)} = \frac{(-1)(-3)}{(1)(-1)} = \frac{3}{-1}
              \equiv 3\cdot 6 \equiv 18 \equiv 4 \pmod 7,\\
  \ell_3(0) &= \frac{(0-1)(0-2)}{(3-1)(3-2)} = \frac{(-1)(-2)}{(2)(1)} = \frac{2}{2} = 1.
\end{align*}
Entonces
\begin{align*}
  S = f(0) &= 6\cdot 3 + 1\cdot 4 + 0\cdot 1 \pmod 7 \\
           &= 18 + 4 + 0 = 22 \equiv \boxed{1} \pmod 7 .
\end{align*}
El \textbf{secreto es $S = 1$}. (Se verifica más abajo que la cuarta sombra $(4,3)$ es
consistente con el mismo polinomio, lo que confirma el resultado.)

\textbf{b) Polinomio.} Buscamos $f(x)=a_0+a_1x+a_2x^2 \bmod 7$ con $f(1)=6$, $f(2)=1$,
$f(3)=0$. Como $a_0 = S = 1$, el sistema queda:
\begin{align*}
  f(1) &= 1 + a_1 + a_2 \equiv 6 &&\Rightarrow\; a_1 + a_2 \equiv 5,\\
  f(2) &= 1 + 2a_1 + 4a_2 \equiv 1 &&\Rightarrow\; 2a_1 + 4a_2 \equiv 0
          \;\Rightarrow\; a_1 + 2a_2 \equiv 0 .
\end{align*}
Restando la primera de la segunda: $a_2 \equiv -5 \equiv 2 \pmod 7$, y entonces
$a_1 \equiv 5 - a_2 = 3 \pmod 7$. El polinomio es
\begin{equation*}
  \boxed{\,f(x) = 1 + 3x + 2x^2 \pmod 7\,}.
\end{equation*}
\textbf{Verificación} de las cuatro sombras:
\begin{align*}
  f(1) &= 1+3+2 = 6, & f(2) &= 1+6+8 = 15 \equiv 1,\\
  f(3) &= 1+9+18 = 28 \equiv 0, & f(4) &= 1+12+32 = 45 \equiv 3 \pmod 7 .
\end{align*}
Coinciden con $C$, incluida la cuarta sombra (que es redundante, como debe ser en un esquema
$(3,4)$).

\textbf{c) Significado vía entropía.} Un esquema umbral $(3,n)$ es válido cuando conocer
\textbf{menos de $3$} sombras \textbf{no aporta información} sobre el secreto $S$ (su entropía
no cambia), y conocer $3$ o más lo determina por completo (entropía $0$):
\begin{equation*}
\begin{aligned}
  H(S \mid S_{i_1})            &= H(S)  && \text{(1 sombra: no sé nada más)},\\
  H(S \mid S_{i_1}, S_{i_2})   &= H(S)  && \text{(2 sombras: tampoco)},\\
  H(S \mid S_{i_1}, S_{i_2}, S_{i_3}) &= 0 && \text{(3 sombras: tengo toda la información)}.
\end{aligned}
\end{equation*}
Es decir: con menos de $T=3$ sombras \textbf{no hay flujo de información} hacia el atacante
($H$ se mantiene en su máximo, $\log_2 7$ bits si $S$ es uniforme en $\mathbb{Z}_7$); recién al
alcanzar el umbral la incertidumbre colapsa a $0$ y se recupera el secreto. Esa es la propiedad
de \emph{seguridad perfecta del umbral} del esquema de Shamir expresada con entropía.

% =====================================================================
\subsection{Ejercicio 4 --- Multiple choice con justificación}
% =====================================================================

\begin{consigna}
En cada caso selecciona la opción correcta. Justificar en todos los casos. Si no se justifica
correctamente, la respuesta quedará invalidada.

\textbf{1.} Si se tiene una lista de control de acceso $\mathrm{ACL}(o) = \{(\text{Manuel}, r),
(\text{Directivos}, w), (*,w)\}$. Asumiendo que Manuel es Directivo,
\begin{enumerate}[label=\alph*.]
  \item Manuel puede leer y escribir en $o$ si la política es first rule.
  \item Manuel puede leer y escribir en $o$ si la política es grant all.
  \item Manuel puede leer y escribir en $o$ si la política es override.
  \item Manuel puede leer y escribir en $o$ si la política es augment.
\end{enumerate}

\textbf{2.} Se conoce como virus polimórfico
\begin{enumerate}[label=\alph*.]
  \item a aquél que infecta distintos tipos de archivos.
  \item a aquél que se modifica cada vez que se replica.
  \item a aquél que se replica muchas veces.
  \item a aquél que se aloja en el sector de arranque o en la memoria.
\end{enumerate}

\textbf{3.} En un sistema de autenticación la Función de Complementación es la que:
\begin{enumerate}[label=\alph*.]
  \item Dada información de autenticación, obtiene información complementaria
  \item Dada información complementaria, obtiene información de autenticación
  \item Dada información de autenticación o información complementaria, las actualiza.
  \item Dada información de autenticación e información complementaria, retorna verdadero o falso
\end{enumerate}
\end{consigna}

\paragraph{Temas.}
Clase 6 --- Políticas y Control de Acceso ($\S$ Conflictos en ACL; $\S$ Derechos por defecto:
override / augment). Clase 9 --- Malware ($\S$ Virus / Worm). Clase 7 --- Autenticación
($\S$ Componentes del sistema de autenticación: tupla $\langle A,C,F,L,S\rangle$).

\paragraph{Qué necesitás saber.}
\textbf{(4.1) Resolución de conflictos de ACL.} Manuel matchea tres entradas: $(\text{Manuel},r)$,
$(\text{Directivos},w)$ (es Directivo) y $(*,w)$ (el comodín/\emph{default}). Las políticas:
\begin{itemize}
  \item \textbf{First rule / first match:} se aplica \emph{solo la primera} entrada que matchea
        $\Rightarrow$ $(\text{Manuel},r)$ $\Rightarrow$ solo \textbf{lee}.
  \item \textbf{Grant all:} se concede un permiso solo si \emph{todas} las entradas que aplican
        lo conceden (intersección). $r \cap w \cap w = \varnothing$ $\Rightarrow$ \textbf{ningún}
        permiso común; ni $r$ ni $w$.
  \item \textbf{Override / best match:} si hay entrada específica se usa esa y se ignora el
        default. La específica para Manuel es $(\text{Manuel},r)$ $\Rightarrow$ solo \textbf{lee}.
  \item \textbf{Augment:} se parte del default y se le \emph{suman} las entradas que aplican
        (unión). Default $(*,w)$ + $(\text{Directivos},w)$ + $(\text{Manuel},r)$
        $= \{r,w\}$ $\Rightarrow$ \textbf{lee y escribe}.
\end{itemize}
\textbf{(4.2) Virus polimórfico.} Un virus polimórfico \textbf{cambia su propio código en cada
replicación} (típicamente re-cifrando su cuerpo con una clave/decodificador distinto) para
evadir la detección por firma del antivirus. Encaja con la familia de \emph{rootkit}/virus que
usan ofuscación/encripción como contramedida (Clase 9, $\S$ Rootkit).
\textbf{(4.3) Función de complementación $F$.} En la tupla $\langle A,C,F,L,S\rangle$,
$F:\,A\to C$ \textbf{deriva} la información complementaria $c$ (lo que guarda el sistema, p.ej.
el hash) a partir de la información de autenticación $a$. No confundir con $L:\,A\times C\to\{0,1\}$,
que \textbf{corrobora} (el login, opción d).

\paragraph{Notas y conexiones.}
En 4.1 la trampa es que \emph{first rule} y \emph{override} dan lo mismo aquí (solo $r$), porque
la entrada específica de Manuel es la primera y a la vez la más específica; ninguna de las dos
le da escritura. Solo \emph{augment} (unión con el default) le da $r$ \emph{y} $w$. En 4.3 es la
distinción \emph{F deriva / L corrobora} que la cátedra marca como pregunta clásica; las
opciones (a) y (d) son justamente $F$ y $L$.

\paragraph{Resolución.}
\begin{itemize}
  \item \textbf{4.1 $\to$ opción (d), augment.} Solo \emph{augment} concede a Manuel
        \textbf{leer y escribir}: parte del default $(*,w)$ y le agrega su $r$ específico,
        resultando $\{r,w\}$. \emph{First rule} da solo $r$ (primera entrada), \emph{grant all}
        no da nada (intersección $r\cap w=\varnothing$) y \emph{override} usa la específica
        $(\text{Manuel},r)$, también solo $r$. Por lo tanto (a), (b) y (c) son falsas.
  \item \textbf{4.2 $\to$ opción (b),} ``se modifica cada vez que se replica''. Esa es la
        definición de polimorfismo (mutar el código en cada copia para evadir firmas). Las
        otras describen: (a) virus multipartito/multi-formato, (c) mera tasa de replicación,
        (d) virus de boot sector/residente en memoria.
  \item \textbf{4.3 $\to$ opción (a).} La función de complementación es $F:\,A\to C$: dada la
        información de autenticación, \emph{deriva} (obtiene) la información complementaria.
        La (d) describe la función de autenticación $L$ (corrobora, retorna verdadero/falso);
        la (c) describe las funciones de selección $S$ (altas/cambios); la (b) sería invertir
        $F$, que justamente debe ser \emph{no invertible}.
\end{itemize}
